\exercisesection{\S~2}

\begin{enumerate}
\def\labelenumi{\arabic{enumi})}
\item
  迹为零的对角矩阵构成 \(\mathfrak{sl}(n,k)\) 的一个嘉当子代数，但当 \(n = 2\) 且 \(k\) 的特征为 2 时除外。
\item
  设 \(e\) 是元素 \(\begin{pmatrix} 0 & 1 \\ 0 & 0 \end{pmatrix}\) ，它属于 \(\mathfrak{sl}(2, \mathbf{C})\) 。证明 \(\mathbf{C}e\) 是 \(\mathfrak{sl}(2, \mathbf{C})\) 的极大幂零李子代数，但不是 \(\mathfrak{sl}(2, \mathbf{C})\) 的嘉当子代数。
\item
  假设 \(k\) 的特征为 0。设 \(\mathfrak{g}\) 是半单李代数。设 \(E\) 是 \(\mathfrak{g}\) 中交换子代数的集合，其元素均在 \(\mathfrak{g}\) 中半单。则 \(\mathfrak{g}\) 的嘉当子代数正是 \(E\) 的极大元。（利用定理 2 与命题 10。）
\end{enumerate}

特别地，\(\mathfrak{g}\) 的诸嘉当子代数之并等于 \(\mathfrak{g}\) 的半单元素组成的集合。

\begin{enumerate}
\def\labelenumi{\arabic{enumi})}
\setcounter{enumi}{3}
\item
  设 \(\mathfrak{g}\) 是具有基 \((x,y,z)\) 的李代数，使得 \([x,y] = y\) ，\([x,z] = z\) ，\([y,z] = 0\) 。设 \(\mathfrak{a}\) 为理想 \(ky + kz\) （\(\mathfrak{g}\) 中的）。则 \(\mathrm{rk}(\mathfrak{a}) = 2\) 且 \(\mathrm{rk}(\mathfrak{g}) = 1\) 。
\item
  假设 \(k\) 的特征为 0。设 \(\mathfrak{g}\) 是李代数，\(\mathfrak{r}\) 是其根基，\(\mathfrak{h}\) 是 \(\mathfrak{g}\) 的一个嘉当子代数。证明
\end{enumerate}

\[
\mathfrak {r} = [ \mathfrak {g}, \mathfrak {r} ] + (\mathfrak {h} \cap \mathfrak {r}).
\]

（注意 \(\mathfrak{h}\) 在 \(\mathfrak{g}/[\mathfrak{g},\mathfrak{r}]\) 中的像包含 \(\mathfrak{r}/[\mathfrak{g},\mathfrak{r}]\) ，即 \(\mathfrak{g}/[\mathfrak{g},\mathfrak{r}]\) 的中心。）

\begin{enumerate}
\def\labelenumi{\arabic{enumi})}
\setcounter{enumi}{5}
\item
  设 \(\mathfrak{g}\) 是李代数，\(\mathfrak{h}\) 是 \(\mathfrak{g}\) 的幂零子代数。若 \(\mathfrak{g}^0 (\mathfrak{h})\) 幂零，则 \(\mathfrak{g}^0 (\mathfrak{h})\) 是 \(\mathfrak{g}\) 的嘉当子代数。
\item
  设 \(\mathfrak{s}\) 是李代数，\(\mathfrak{a}\) 是 \(\mathfrak{s}\) 的嘉当子代数，V 是 \(\mathfrak{s}\) -模。设 \(\mathfrak{g} = \mathfrak{s} \times V\) 是 \(\mathfrak{s}\) 与 V 的半直积。证明 \(\mathfrak{a} \times V^0(\mathfrak{a})\) 是 \(\mathfrak{g}\) 的嘉当子代数。
\item
  假设 \(k\) 的特征为 \(p > 0\) 。用 \(\mathfrak{s}\) 表示具有基 \(\{x, y\}\) 且满足 \([x, y] = y\) 的李代数。设 \(V\) 是 \(k\) -向量空间，具有基 \(\{e_i\}_{i \in \mathbf{Z}/p\mathbf{Z}}\) 。
\end{enumerate}

\begin{enumerate}
\def\labelenumi{\alph{enumi})}
\item
  证明 V 有唯一的 \(\mathfrak{s}\) -模结构，使得 \(xe_{i} = ie_{i}\) 且 \(ye_{i} = e_{i + 1}\) 对一切 \(i\) 成立。此 \(\mathfrak{s}\) -模是单模。
\item
  设 \(\mathfrak{g} = \mathfrak{s} \times \mathrm{V}\) 是 \(\mathfrak{s}\) 与 \(\mathrm{V}\) 的半直积。证明 \(\mathfrak{g}\) 是秩为 1 的可解代数，其导代数不幂零。
\item
  \(\mathfrak{g}\) 的一个元素是正则元，当且仅当它在 \(\mathfrak{s}\) 上的投影形如 \(ax + by\) ，其中 \(ab \neq 0\) 。
\item
  我们有 \(\mathrm{V}^{0}(x+y)=0\) 且 \(\mathrm{V}^{0}(x)=ke_{0}\) 。推出（参见习题 6）g 有维数为 1 的嘉当子代数（例如由 \(x+y\) 生成的那个），以及维数为 2 的嘉当子代数（例如由 x 与 \(e_{0}\) 生成的那个）。
\end{enumerate}

\begin{enumerate}
\def\labelenumi{\arabic{enumi})}
\setcounter{enumi}{8}
\item
  设 g 是具有基 \((x,y)\) 且满足 \([x,y]=y\) 的李代数。设 k=ky，并用 \(\varphi:g\to g/k\) 表示典则态射。g/k 的元素 0 在 g/k 中正则，但它不是 g 的正则元在 \(\varphi\) 下的像。
\item
  假设 \(k\) 无限。设 \(\mathfrak{g}\) 是李代数。证明下列性质等价：
\end{enumerate}

\begin{enumerate}
\def\labelenumi{(\roman{enumi})}
\item
  \(\operatorname{rk}(\mathfrak{g}) = \dim (\mathfrak{g})\)。
\item
  \(\mathfrak{g}\) 是幂零的。
\item
  g 只有有限多个维数为 rk(g) 的嘉当子代数。
\item
  g 只有一个嘉当子代数。
\end{enumerate}

\begin{enumerate}
\def\labelenumi{\arabic{enumi})}
\setcounter{enumi}{10}
\tightlist
\item
  设 \(\mathfrak{h}\) 是一个 \(\neq 0\) 的交换李代数，\(P\) 是 \(\mathfrak{h}^*\) 的包含 0 的有限子集。证明存在一个李代数 \(\mathfrak{g}\) ，以 \(\mathfrak{h}\) 为嘉当子代数，且使得 \(\mathfrak{h}\) 在 \(\mathfrak{g}\) 中的权集为 \(P\) 。（把 \(\mathfrak{g}\) 构造为 \(\mathfrak{h}\) 与 \(\mathfrak{h}\) -模 \(V\) 的半直积，其中 V 是对应于 \(P - \{0\}\) 中诸元素的一维模的直和，参见习题 7。）
\end{enumerate}

\(x\) 是 \(\mathfrak{h}\) 中的元素，满足 \(\mathfrak{h} = \mathfrak{g}^0(x)\) ，当且仅当 \(x\) 不与 \(\mathrm{P} - \{0\}\) 中任何元素正交。

\begin{enumerate}
\def\labelenumi{\arabic{enumi})}
\setcounter{enumi}{11}
\tightlist
\item
  假设 \(k\) 有限。构造一个李代数 \(\mathfrak{g}\) 的例子，它有一个嘉当子代数 \(\mathfrak{h}\) ，其中不存在元素 \(x\) 使得 \(\mathfrak{h} = \mathfrak{g}^0(x)\) 。（利用前一习题，并取 \(\mathrm{P} = \mathfrak{h}^*\) 。）
\end{enumerate}

¶ 13) 假设 k 有限。用 \(k'\) 表示 k 的一个无限扩张。设 g 是 k 上的李代数。g 的秩记作 \(\operatorname{rk}(\mathfrak{g})\) ，定义为 \(k'\) -李代数 \(g' = g \otimes_{k} k'\) 的秩；g 的元素称为正则的，若它在 \(g'\) 中正则；这些定义不依赖于 \(k'\) 的选取。证明：若

\[
\operatorname{Card} (k) \geq \dim \mathfrak {g} - \operatorname{rk} (\mathfrak {g}),
\]

\(\mathfrak{g}\) 含有正则元（从而也含有嘉当子代数）。

（利用下述结果：若 \(a\) 是 \(k[\mathrm{X}_1, \ldots, \mathrm{X}_n]\) 的非零齐次元素，且 \(\operatorname{Card}(k) \geq \deg(a)\) ，则存在 \(x \in k^n\) 使得 \(a(x) \neq 0\) 。）

\begin{enumerate}
\def\labelenumi{\arabic{enumi})}
\setcounter{enumi}{13}
\tightlist
\item
  假设 \(k\) 的特征为零。设 V 是有限维 \(k\) -向量空间，\(\mathfrak{g}\) 是 \(\mathfrak{gl}(V)\) 的李子代数，\(\mathfrak{h}\) 是 \(\mathfrak{g}\) 的嘉当子代数，\(\mathfrak{n}_{\mathrm{V}}\) 是 \(\mathfrak{g}\) -模 V 的最大幂零理想（第一章 §4 第 3 节 定义 2）。
\end{enumerate}

证明 h 的元素幂零，当且仅当它属于 \(n_{V}\) 。（化归到 g-模 V 半单的情形，并利用定理 2 的推论 3。）

¶ 15) 假设 \(k\) 无限。设 \(\mathfrak{g}\) 是李代数，\(\mathcal{C}_{\infty}\mathfrak{g}\) 是 \(\mathfrak{g}\) 的升中心列之并，\(x\) 是 \(\mathfrak{g}\) 的元素。证明下列性质等价：

\begin{enumerate}
\def\labelenumi{(\roman{enumi})}
\item
  \(x\) 属于 \(\mathfrak{g}\) 的每个嘉当子代数。
\item
  \(x\in \mathfrak{g}^0 (y)\) 对一切 \(y\in \mathfrak{g}\) 成立（即 \(x\in \mathfrak{g}^0 (\mathfrak{g})\) ）。
\item
  \(x \in C_{\infty}g.\)
\end{enumerate}

（蕴涵关系 (iii) \(\Longrightarrow\) (ii) \(\Longrightarrow\) (i) 是直接的。为证 (i) \(\Longrightarrow\) (ii)，注意 (i) 等价于说 \(x\in \mathfrak{g}^0 (y)\) 对每个正则元 \(y\) （\(\mathfrak{g}\) 中的）成立，并利用正则元在 Zariski 拓扑下于 \(\mathfrak{g}\) 中稠密这一事实。为证 (ii) \(\Longrightarrow\) (iii)，注意 \(\mathfrak{n} = \mathfrak{g}^0 (\mathfrak{g})\) 在 \(\mathfrak{g}\) 下稳定（§1 习题 5），并对 \(\mathfrak{g}\) -模 \(\mathfrak{n}\) 应用 Engel 定理；推出 \(\mathfrak{n}\) 含于 \(\mathcal{L}_{\infty}\mathfrak{g}\) 。）

¶ 16) 设 g 是可解复李代数，h 是 g 的嘉当子代数，\(\mathfrak{g} = \bigoplus \mathfrak{g}^{\lambda}(\mathfrak{h})\) 是 g 的相应准素子空间分解，其中 \(\mathfrak{g}^{0}(\mathfrak{h}) = \mathfrak{h}\) 。

\begin{enumerate}
\def\labelenumi{\alph{enumi})}
\item
  证明线性形式在 \(\mathfrak{h}\) 上的限制——这些线性形式即第三章 §9 习题 17 c) 中称为 \(\mathfrak{g}\) 的根者——是 \(\mathfrak{h}\) 在 \(\mathfrak{g}\) 中的权，即诸 \(\lambda\) ，使得 \(\mathfrak{g}^{\lambda}(\mathfrak{h}) \neq 0\) ；推出这样的 \(\lambda\) 在 \(\mathfrak{h} \cap \mathcal{D}\mathfrak{g}\) 上为零。
\item
  设 \((x,y)\mapsto [x,y]'\) 是从 \(\mathfrak{g}\times \mathfrak{g}\) 到 \(\mathfrak{g}\) 的交错双线性映射，具有下列性质：
\end{enumerate}

\begin{enumerate}
\def\labelenumi{(\roman{enumi})}
\item
  若 \(x \in \mathfrak{g}^{\lambda}(\mathfrak{h}), y \in \mathfrak{g}^{\mu}(\mathfrak{h})\) ，其中 \(\lambda \neq 0, \mu \neq 0\) ，则 \([x, y]' = [x, y]\) ；
\item
  若 \(x \in \mathfrak{g}^{0}(\mathfrak{h}), y \in \mathfrak{g}^{\mu}(\mathfrak{h})\) ，则 \([x, y]' = [x, y] - \mu(x)y\) 。
\end{enumerate}

证明这在 \(\mathfrak{g}\) 上给出一个新的李代数结构（利用 \(a\) ）。把它记作 \(\mathfrak{g}'\) 。

\begin{enumerate}
\def\labelenumi{\alph{enumi})}
\setcounter{enumi}{2}
\tightlist
\item
  证明：若 \(x \in \mathfrak{g}^{\lambda}(\mathfrak{h})\) ，则映射 \(ad'x : y \mapsto [x, y]'\) 幂零。推出 \(g'\) 幂零（把第一章 §4 习题 11 应用于 \(ad'x\) 的集合 E，其中 x 属于诸 \(\mathfrak{g}^{\lambda}(\mathfrak{h})\) 之并）。
\end{enumerate}

